% answers_short.tex —— 2010 年·普通高等学校招生考试·江西卷·理科 · 简短答案速查
% 由 tools/make_short_answers.py 自 answers.tex 的 \gksolhead/\gkanswer 提取，
% 并与 2010/2010_数学真题_江西理_答案要点.md 逐题比对后生成，勿手改（改源 MD 后重跑）。

\begin{gknotice}
\textbf{使用说明}：本节只给最终答案便于快速核对，完整推导见解析版。答案为编写者独立推导，\textbf{非官方标准答案}。
\end{gknotice}

\gkpart{选择题}
\begin{center}\begin{tabular}{ccccccccccccc}
\toprule
题号 & 1 & 2 & 3 & 4 & 5 & 6 & 7 & 8 & 9 & 10 & 11 & 12 \\
\midrule
答案 & D & C & A & B & C & B & D & A & C & D & B & D \\
\bottomrule
\end{tabular}\end{center}

\gkpart{填空题}
\begin{itemize}[leftmargin=2.6em]
  \item[\textbf{13.}] $\sqrt3$。
  \item[\textbf{14.}] $1080$。
  \item[\textbf{15.}] $2$。
  \item[\textbf{16.}] $S_1>S_2>S_3$。
\end{itemize}

\gkpart{解答题}
\begin{itemize}[leftmargin=2.6em]
  \item[\textbf{17.}] （1）$\left[0,\dfrac{1+\sqrt2}{2}\right]$；（2）$m=-2$。
  \item[\textbf{18.}] （1）$P(\xi=1)=\dfrac13$，$P(\xi=3)=\dfrac16$，$P(\xi=4)=\dfrac16$，$P(\xi=6)=\dfrac13$；（2）$E\xi=\dfrac72$。
  \item[\textbf{19.}] （1）增区间 $(0,\sqrt2)$，减区间 $(\sqrt2,2)$；（2）$a=\dfrac12$。
  \item[\textbf{20.}] （1）$\dfrac{2\sqrt{15}}5$；（2）$\dfrac{2\sqrt5}5$。
  \item[\textbf{21.}] （1）$\dfrac{\sqrt2}{2}$；（2）$C_1:\dfrac{3x^2}{16}+\dfrac{y^2}4=1$，$C_2:x^2+2y=4$。
  \item[\textbf{22.}] （1）$b=5a$，$c=7a$ 即可；（2）取 $a=m^2-2mn-n^2$，$b=m^2+n^2$，$c=m^2+2mn-n^2$（$m,n$ 为正整数且 $(m^2-n^2)^2>12m^2n^2$），得无穷多个互不相似的三角形。
\end{itemize}
